\section{A fast-forwarding theory explorer}
\label{sec:lift}
\label{subsec:ff-nointerp}
\label{subsec:perf}

This section presents a new
  \textit{fast-forwarding} theory exploration
  algorithm, which has two key applications.
First, as \autoref{subsec:ffoverview} showed,
  it enables rewrite rule inference
  for domains where writing an interpreter
  is prohibitively difficult.
Second,
  it mitigates performance problems that arise
  when resource limits interact poorly with
  the kinds of rules in a ruleset.

The \textit{kinds} of rules that comprise a ruleset
  significantly affect performance,
  even in efficient \eqsat-driven systems.
Since reaching saturation
  in an \egraph is rare in practice,
  iteration and/or node limits are used
  to ensure termination.
As a result,
  two rulesets can have vastly different
  performance
  even if they are equivalent under derivability.

The key motivation behind
  the fast-forwarding algorithm is that
  the ``right'' set of rules
  can help \textit{fast-forward} equality saturation by
  skipping intermediate derivations.
Skipping intermediate derivations has two benefits:
  first, it requires fewer iterations to prove
  a target equivalence;
  second, it often
  reduces the number of intermediate terms in the \egraph,
  reducing unhelpful rewriting on these terms.
Determining the ``right'' rules requires domain knowledge
  and depends on the application.
To that end, we assume that the user can provide
  a set of \textit{allowed} ($\mathcal{A}$) and
  \textit{forbidden} ($\mathcal{F}$) operators.
We then say that
  if a pattern $p$ contains \textit{any} operator $o \in \mathcal{F}$,
  then $p$ is forbidden.
If \textit{all} operators in $p$ are allowed,
  then the pattern is allowed.
Since a rule is simply a pair
  of patterns,
  these definitions extend to rules.

For example, in the case of the trigonometric
  rule synthesis task in \autoref{subsec:ffoverview},
  the allowed operators are \texttt{sin}, \texttt{cos}, \texttt{tan},
  \texttt{PI}, \texttt{+, -, $\times$}, etc., and
  the forbidden operators are
  \texttt{cis} and \texttt{I} because we wanted \enumo to synthesize
  rewrite rules over the trigonometric domain only.
Recall that this task also required an additional set of
  \textit{exploratory} ($\mathcal{E}$) rewrite rules that relate
  terms with allowed operators to terms with \emph{other} operators.
Crucially, these ``other'' operators can be either allowed or
  forbidden.
The intuition behind $\mathcal{E}$
  is that it helps \textit{explore}
  new equivalences between allowed terms in the \egraph by
  applying a \textit{known} set of rewrites over terms containing the
  \emph{other} operators
  (shown by \texttt{explore} in \autoref{subsec:ffoverview}).

\begin{figure}
\begin{lstlisting}
def fast_forward_naive($\mathcal{W},\, \mathcal{R}$):
  $\mathcal{G}$ = $\mathcal{W}$.to_egraph() # convert workload to e-graph
  $\mathcal{G'}$ = $\mathcal{G}$.compress($\mathcal{R}$) # run equality saturation with all rules
  # any two terms from the unions are potential candidates:
  candidates = candidates_by_diff($\mathcal{G}$, $\mathcal{G'}$)
  return candidates.minimize($\mathcal{R}$) # minimize candidates
\end{lstlisting}
  \caption{A naive algorithm for fast-forwarding theory exploration which
  applies \eqsat to terms represented by $\mathcal{W}$ using
  \T{compress}.}
\label{fig:ff-naive}
\end{figure}

\autoref{fig:ff-naive} shows a naive implementation
  of fast-forwarding
  using the set of core \enumo operators
  from \autoref{sec:slide}.
The process consists of
  applying \T{eqsat} to a workload representing
  \textit{allowed} terms
  using a ruleset, $\mathcal{R}$, which contains both
  allowed and forbidden rules.
First, it creates an \eggl from the terms obtained by
  evaluating the workload.
Then, it \textit{shrinks} the \egraph using the
  \T{compress} operator.
\T{compress} is an \eqsat \textit{strategy} (\autoref{fig:ruleops}) that prevents
  the \egraph from getting intractably large.
It applies the rewrites on a duplicate of the original
  \egraph and only copies the
  equivalences back.
It adds no new \enodes or \eclasses to the original \egraph.
The next step in the algorithm extracts candidates from $\mathcal{G}$
  based on the equalities discovered in $\mathcal{G'}$
  using a cost function that penalizes forbidden operators.
Finally, it minimizes the resulting ruleset as explained
  in \autoref{sec:slide}.
Notice that this naive algorithm
  simply performs a single phase of \T{compress} with
  \textit{all} the rules.

\begin{figure}
\begin{lstlisting}
def fast_forward($\mathcal{W},\, \mathcal{R},\, \mathcal{E}$):
  $\mathcal{G}$ = $\mathcal{W}$.to_egraph() $\label{ln:5}$ # convert workload to egraph
  allowed = $\{r ~ \in ~ \mathcal{R} ~ | ~ \forall ~ o ~\in~ (ops(r.lhs) ~ \cup ~ ops(r.rhs)), ~ o ~\in~ \mathcal{A} \}$  $\label{ln:6}$
  $\mathcal{G'}$ = $\mathcal{G}$.compress(allowed) $\label{ln:7}$ # compress the egraph with allowed rules
  $\mathcal{G''}$ = $\mathcal{G'}$.eqsat($\mathcal{E}$) $\label{ln:11}$ # grow the egraph with exploratory rules
  $\mathcal{G'''}$ = $\mathcal{G''}$.compress($\mathcal{R}$) $\label{ln:12}$ # compress to find equalities with all of $\mathcal{R}$
  candidates = candidates_by_diff($\mathcal{G'}$, $\mathcal{G'''}$) $\,$# extract learned rules with no ops in $\mathcal{F}$
  return candidates.minimize(allowed) # minimize candidates
\end{lstlisting}
\caption{A practical fast-forwarding theory exploration algorithm that
  approximates the naive version. \textit{ops} is a helper function that returns
  all the operators in a term.}
\label{fig:ff-actual}
\end{figure}

\paragraph*{A practical algorithm.}
Unfortunately, the naive algorithm in \autoref{fig:ff-naive}
  does not find useful
  rules in practice for two reasons.
First, it does not scale:
  for large workloads, the \egraph blows up and
  resource limits (e.g., timeout, iterations) are exhausted
  before useful equivalences emerge.
Second, exploring in a breadth-first manner
  prevents finding interesting fast-forwarding opportunities,
  which only occur after several rounds of \eqsat.

Instead, we propose a more practical, approximate algorithm
  that applies \eqsat
  in a more strategic way by leveraging a user's domain
  knowledge in the form of $\mathcal{E}$,
  $\mathcal{F}$, and $\mathcal{A}$.
The algorithm in \autoref{fig:ff-actual}
  \textit{selectively} grows and compresses
  the \egraph using the rules provided by the user.
The algorithm first creates an \eggl
  from the terms represented by $\mathcal{W}$,
  then compresses the \egraph with allowed rules
  (lines \ref{ln:6}--\ref{ln:7}).
This step  \textit{shrinks} the \egraph with known
  equivalences.
In the next step,
  the algorithm \textit{grows} the \egraph with
  $\mathcal{E}$ (line \ref{ln:11}).
Crucially, this step does \emph{not} use \T{compress};
  it performs simple \T{eqsat} that introduces
  new terms and equivalences in the \egraph.
The final \eqsat step applies another round of compression
  using all the rules in $\mathcal{R}$, discovering new
  fast-forwarded rules.
The minimization step in this algorithm
  uses the \T{allowed} rules instead of the entire
  ruleset in order to avoid forbidden operators
  in the minimized ruleset.
The rest of the algorithm is similar to \autoref{fig:ff-naive}.

\subsection{Comparing different scheduling strategies}
\label{subsec:cmp}
The key idea in our fast-forwarding algorithm is to
  perform \eqsat in phases,
  using subsets of $\mathcal{R}$ to selectively
  grow and compress the \egraph.
To understand how this affects performance,
  we ran an experiment to evaluate and compare
  the difference between using
  \C{eqsat} and \C{compress}
  in \autoref{fig:ff-actual}.
We used a workload of 287 terms
  from the domain of trigonometric operators
  ($\sin$, $\cos$, $\tan$, $\pi$, $\pi/2$, etc.).
\autoref{table:ff-cmp} shows the
  results of the comparison.
The first two rows, which use \C{eqsat} in all three phases, do not terminate
  within 20 minutes.
These treatments demonstrate the importance of \C{compress}, which does not
  allow the \egraph to grow.
The next two rows use \C{compress} in all three phases.
The third row does not split up the rules in $\mathcal{R}$ and simply runs
  \C{compress($\mathcal{R}$)} three times.
The fourth row compresses with the allowed rules ($\mathcal{A}$) in
  Phase 1, the exploratory rules ($\mathcal{E}$)
  in Phase 2, and all rules ($\mathcal{R}$) in Phase 3.
Both of these treatments finish within seconds, but do not find any new rules.
The approach in the last row, which corresponds to the actual \ff algorithm
  described in \autoref{fig:ff-actual}, finds 4 useful trigonometric
  identities in about 3 minutes.
This experiment demonstrates the importance of using \C{eqsat} and \C{compress}
  together to strategically grow and compress the \egraph.

\begin{table}[h]
  \footnotesize
\begin{tabular}{llllc}
  Phase 1 & Phase 2 & Phase 3 & Time (s) & \# Rules with Trig Operators \\ \cline{1-5}
eqsat ($\mathcal{R}$) & eqsat ($\mathcal{R}$) & eqsat ($\mathcal{R}$) & Timeout & - \\
eqsat ($\mathcal{A}$) & eqsat ($\mathcal{E}$) & eqsat ($\mathcal{R}$) & Timeout & - \\
compress ($\mathcal{R}$) & compress ($\mathcal{R}$) & compress ($\mathcal{R}$) & 18.66 & 0 \\
compress ($\mathcal{A}$) & compress ($\mathcal{E}$) & compress ($\mathcal{R}$) & 9.42 & 0 \\
compress ($\mathcal{A}$) & eqsat ($\mathcal{E}$) & compress ($\mathcal{R}$) & 175.78 & 4 \\
\end{tabular}%
\caption{Comparing \texttt{compress} and \texttt{eqsat} with different
    subsets of $\mathcal{R}$ for the three phases of \autoref{fig:ff-actual}.
    $\mathcal{A}$ is the allowed rules of $\mathcal{R}$ and
    $\mathcal{E}$ is the exploratory rules of $\mathcal{R}$.
    The last row corresponds to \autoref{fig:ff-actual};
    it is the only treatment that finds any trigonometric rules.}
\label{table:ff-cmp}
\end{table}

